\section{Real-Time Windowing}
\label{sec:windowing}

The accelerometer never stops: at \SI{64}{Hz} it hands Cadence a fresh triple of
$(a_x,a_y,a_z)$ samples roughly every \SI{16}{ms}, indefinitely. A freeze, by
contrast, is an \emph{event} with duration --- a few seconds of \SIrange{3}{8}{Hz}
trembling. To bridge that gap, the firmware does not classify single samples. It
slices the unending stream into short, overlapping \emph{windows}, computes one
spectral verdict per window, and slides forward. This section is about the geometry
of that slicing: how wide each window is, how far it steps, what that buys us, and
what it quietly costs.

\begin{intuition}
You cannot hear a chord from a single piano key struck for one millisecond --- you
need to listen for a beat or two to tell which notes are sounding. The freeze
signature is a ``chord'' in the \SIrange{3}{8}{Hz} band. Cadence listens in
\SI{4}{s} chunks, and starts a fresh chunk every \SI{2}{s}, so a new opinion arrives
twice a second while each opinion still rests on enough signal to be trustworthy.
\end{intuition}

\subsection{The streaming geometry}

Two integers fix everything. The \textbf{window length} is $N=256$ samples; at
$f_s=\SI{64}{Hz}$ that is exactly
\[
  T_{\text{win}} \;=\; \frac{N}{f_s} \;=\; \frac{256}{64} \;=\; \SI{4.0}{s}.
\]
The \textbf{hop} (how far the window advances between decisions) is $H=128$ samples,
i.e.\ \SI{2.0}{s}. Because $H = N/2$, consecutive windows share exactly half their
samples --- a \textbf{50\% overlap} --- and a fresh window closes, and therefore a
fresh freeze decision is emitted, every
\[
  T_{\text{hop}} \;=\; \frac{H}{f_s} \;=\; \frac{128}{64} \;=\; \SI{2.0}{s}.
\]

\begin{keyidea}
Window length sets \emph{how much signal each decision sees}; hop sets \emph{how
often a decision is made}. They are independent knobs. Cadence deliberately decouples
them: a long \SI{4}{s} window for spectral quality, but a short \SI{2}{s} hop so the
monitor still feels responsive. Overlap is simply the by-product, $1-H/N = 50\%$.
\end{keyidea}

Figure~\ref{fig:win-timeline} draws the geometry directly. Each bar is one \SI{4}{s}
analysis window; each window is shifted \SI{2}{s} to the right of the one before, so
its leading half is brand-new signal and its trailing half was already seen by its
predecessor. Reading down any vertical line tells you how many windows a given
instant of the stream participates in --- in steady state, every sample is examined
by exactly two windows.

\begin{figure}[H]
\centering
\resizebox{0.95\textwidth}{!}{%
\begin{tikzpicture}[font=\small]
  % ---- parameters: 1 s = 0.9 cm on the x-axis ----
  \def\sx{0.9}     % cm per second
  \def\wid{4}      % window width in seconds
  \def\hop{2}      % hop in seconds

  % ---- timeline axis ----
  \draw[->,thick,ink] (0,0) -- ({16.8*\sx},0) node[right]{\,$t$\,(s)};
  \foreach \t in {0,2,4,6,8,10,12,14,16}{
    \draw[ink] ({\t*\sx},0.06) -- ({\t*\sx},-0.06);
    \node[below,ink,font=\footnotesize] at ({\t*\sx},-0.10) {\t};
  }

  % ---- five overlapping windows, stacked upward ----
  \def\wha{0.46}   % half-height of each window bar
  \def\gap{1.18}   % vertical centre-to-centre spacing
  \foreach \k [count=\i from 1] in {0,2,4,6,8}{
    \pgfmathsetmacro{\yc}{1.05 + \k/2*\gap}     % vertical centre of window i
    \pgfmathsetmacro{\xl}{\k*\sx}               % left edge (s -> cm)
    \pgfmathsetmacro{\xr}{(\k+\wid)*\sx}        % right edge
    \pgfmathsetmacro{\xm}{(\k+\hop)*\sx}        % overlap boundary (new vs reused)
    % reused (overlap) half — pale
    \filldraw[fill=palenavy,draw=navy,line width=0.5pt,rounded corners=1pt]
      (\xl,{\yc-\wha}) rectangle (\xm,{\yc+\wha});
    % fresh half — teal tint
    \filldraw[fill=paleteal,draw=teal,line width=0.5pt,rounded corners=1pt]
      (\xm,{\yc-\wha}) rectangle (\xr,{\yc+\wha});
    \node[ink,font=\footnotesize] at ({(\xl+\xr)/2},\yc)
      {window $w_{\i}$ \;[\SI{4}{s}, $N{=}256$]};
    % a decision is emitted when the window closes (its right edge)
    \draw[freeze,line width=0.7pt,-{Stealth[length=4pt]}]
      (\xr,{\yc-\wha}) -- ({\xr},0.18);
    \node[freeze,font=\scriptsize,fill=white,inner sep=0.5pt]
      at ({\xr+0.0},0.40) {decision};
  }

  % ---- hop bracket between w1 and w2 left edges ----
  \draw[decorate,decoration={brace,amplitude=4pt,mirror},steel,line width=0.6pt]
    (0,0.50) -- ({\hop*\sx},0.50);
  \node[steel,font=\footnotesize,fill=white,inner sep=1pt]
    at ({\hop*\sx/2},0.66) {hop $H=128$ (\SI{2}{s})};

  % ---- legend (placed clear of the windows, lower-right) ----
  \begin{scope}[shift={({12.6*\sx},{1.05+1.0*\gap})}]
    \filldraw[fill=paleteal,draw=teal,line width=0.5pt] (0,0) rectangle (0.34,0.30);
    \node[right,ink,font=\footnotesize] at (0.40,0.15) {fresh \SI{2}{s}};
    \node[right,ink,font=\footnotesize] at (0.40,-0.20) {(new signal)};
    \filldraw[fill=palenavy,draw=navy,line width=0.5pt] (0,-0.85) rectangle (0.34,-0.55);
    \node[right,ink,font=\footnotesize] at (0.40,-0.70) {reused \SI{2}{s}};
    \node[right,ink,font=\footnotesize] at (0.40,-1.05) {(overlap)};
  \end{scope}
\end{tikzpicture}}
\caption{Sliding-window timeline. Each \SI{4}{s} ($N{=}256$) analysis window steps
\SI{2}{s} ($H{=}128$) to the right of the one before, so half of every window is new
signal (teal) and half is carried over from its predecessor (navy). A freeze decision
(red arrow) is emitted each time a window closes --- a fresh verdict every \SI{2}{s}.
In steady state every sample is seen by exactly two windows.}
\label{fig:win-timeline}
\end{figure}

\begin{lstlisting}[style=py]
# Streaming geometry, in software (fog.streaming, schematic)
N, H = 256, 128                 # window = 4 s, hop = 2 s  @ 64 Hz
ring = RingBuffer(capacity=N)   # holds the most recent 4 s of samples
for sample in accel_stream():   # 64 new samples per second
    ring.push(sample)
    if ring.advanced_by(H):     # a full 2 s of fresh data has arrived
        window = ring.last(N)   # newest 256 samples (50% overlap with previous)
        verdict = classify(window)   # gate -> freeze index -> debounce
        emit(verdict)           # one decision every 2 s
\end{lstlisting}

\subsection{Why \SI{4}{s}? Spectral resolution versus latency}

The window length is not a free choice --- it is dictated by the physics of the
signal we are trying to resolve. The freeze band, \SIrange{3}{8}{Hz}, must be
separated cleanly from the locomotor band, \SIrange{0.5}{3}{Hz}, sitting right
beside it. A Welch periodogram over $N$ samples places its frequency estimates on a
grid whose spacing is
\[
  \Delta f \;=\; \frac{f_s}{N} \;=\; \frac{64}{256} \;=\; \SI{0.25}{Hz}.
\]
At \SI{0.25}{Hz} per bin the \SIrange{3}{8}{Hz} freeze band is covered by roughly
twenty bins and the \SIrange{0.5}{3}{Hz} walking band by ten --- ample resolution to
add up each band's power independently and form the freeze index as their ratio.

\begin{mathsbox}
$\Delta f = f_s/N$ is the single most important line in this section. The bin spacing
is \emph{inversely} proportional to window length: double the samples, halve the bin
width, sharpen the spectrum. A short \SI{1}{s} window ($N=64$) would give
$\Delta f = \SI{1}{Hz}$ --- so coarse that the \SIrange{3}{8}{Hz} freeze band would
collapse into a handful of fat bins and start to smear into the walking band beneath
it. The \SI{4}{s} window is the smallest that resolves the two bands cleanly.
\end{mathsbox}

There is a price. A longer window means each decision is an \emph{average} over more
time, so it reacts more slowly and blurs together whatever happened across its
\SI{4}{s} span. This is the classic resolution--latency trade-off, summarised in
Table~\ref{tab:tradeoff}.

\begin{table}[H]
\centering
\caption{The window-length trade-off. Cadence sits at $N=256$ (bold row): the
shortest window that still resolves the freeze band, paired with a half-window hop so
decisions still arrive twice a second.}
\label{tab:tradeoff}
\small
\begin{tabular}{@{}lcccl@{}}
\toprule
Window $N$ & $T_{\text{win}}$ & $\Delta f = f_s/N$ & Freeze-band bins & Character \\
\midrule
64  & \SI{1}{s} & \SI{1.0}{Hz}  & $\sim$5  & fast, but bands smear together \\
128 & \SI{2}{s} & \SI{0.5}{Hz}  & $\sim$10 & marginal resolution \\
\textbf{256} & \textbf{\SI{4}{s}} & \textbf{\SI{0.25}{Hz}} & \textbf{$\sim$20} & \textbf{\textcolor{amber}{Cadence: clean bands, responsive}} \\
512 & \SI{8}{s} & \SI{0.125}{Hz} & $\sim$40 & sharp, but sluggish and over-smoothed \\
\bottomrule
\end{tabular}
\end{table}

\begin{intuition}
Resolution and latency pull in opposite directions. A wide window is a
\emph{telephoto lens} on frequency --- it sees the \SIrange{3}{8}{Hz} band in sharp
detail but takes a long exposure, so fast changes blur. A narrow window is a
\emph{snapshot} --- quick, but the frequency picture is grainy. Because a freeze
lasts \emph{seconds}, a \SI{4}{s} exposure is well matched to the event itself: long
enough to be sharp, short enough not to miss the episode.
\end{intuition}

The hop rescues the responsiveness that the long window would otherwise cost us.
Even though each window spans \SI{4}{s}, stepping by only \SI{2}{s} means the latest
decision is never more than \SI{2}{s} old. Overlap is therefore not wasteful
redundancy --- it is what lets a slow-exposure window still produce a fast-updating
verdict. Figure~\ref{fig:capture1} shows this on a real worksheet capture: the freeze
index does not vary smoothly but jumps in a \emph{staircase}, one tread every
\SI{2}{s}, each tread being the verdict of one closed window. The index climbs hard
($>20$) only while the foot is trembling in place during the freeze phase, and the
\SI{2}{s} cadence of the steps is exactly the hop drawn in
Figure~\ref{fig:win-timeline}.

\begin{figure}[H]
\centering
\includegraphics[width=0.86\linewidth]{capture1_trace}
\caption{Real worksheet capture (\code{capture1.csv}, \SI{64}{Hz}, \SI{41.4}{s}).
\emph{Top:} accelerometer magnitude $\norm{a}$, with the freeze phase showing the
tell-tale large-amplitude \SIrange{3}{8}{Hz} trembling. \emph{Bottom:} the freeze
index. Note that it updates as a \emph{staircase} --- one new tread every \SI{2}{s}
hop, each tread the verdict of one \SI{4}{s} window --- and spikes above 20 only
during the freeze, where the worksheet records a mean FI of $6.89$ and a max of
$21.35$.}
\label{fig:capture1}
\end{figure}

\subsection{The ring buffer}

Holding ``the most recent \SI{4}{s}'' on a SAMD21 with kilobytes of RAM rules out
ever copying the stream. Cadence uses a \textbf{ring (circular) buffer}: a fixed
array of $N=256$ slots and a write pointer that wraps around to the start when it
reaches the end. Each new sample overwrites the oldest one in $O(1)$ time, with no
shifting and no allocation. A counter tracks how many fresh samples have landed
since the last decision; the moment it reaches the hop $H=128$, the firmware reads
off the newest 256 slots as the current window and resets the counter.

\begin{figure}[H]
\centering
\resizebox{0.80\textwidth}{!}{%
\begin{tikzpicture}[font=\small]
  \def\R{2.05}
  % ring of slots
  \draw[line width=10pt,rule] (0,0) circle (\R);
  % "fresh" arc (newest 128) in teal, "older" arc in navy
  \draw[line width=10pt,teal] (0,0) ++(90:\R) arc (90:-90:\R);
  \draw[line width=10pt,navy!70] (0,0) ++(90:\R) arc (90:270:\R);
  % write pointer
  \draw[-{Stealth[length=6pt]},freeze,line width=1.3pt] (0,0) -- (88:\R*0.92);
  \node[freeze,font=\footnotesize] at (55:\R*0.50) {write ptr};
  % wrap arrow
  \draw[-{Stealth[length=5pt]},ink,line width=0.8pt]
    (102:\R*1.20) arc (102:78:\R*1.20);
  \node[ink,font=\scriptsize,align=center] at (90:\R*1.42)
    {newest sample\\overwrites oldest};
  % leader lines + labels for the two halves (placed well clear of the 10pt ring)
  \draw[teal,line width=0.5pt] (0:\R*1.08) -- (0:\R*1.30);
  \node[teal,font=\footnotesize,align=left,anchor=west] at (0:\R*1.34)
    {freshest \SI{2}{s}\\(128 slots)};
  \draw[navy,line width=0.5pt] (180:\R*1.08) -- (180:\R*1.30);
  \node[navy,font=\footnotesize,align=right,anchor=east] at (180:\R*1.34)
    {previous \SI{2}{s}\\(reused next hop)};
  % centre annotation
  \node[ink,font=\footnotesize,align=center] at (0,-0.05)
    {256 slots\\$=\SI{4}{s}$\\ at \SI{64}{Hz}};
\end{tikzpicture}}
\caption{The ring buffer holds exactly $N=256$ samples (\SI{4}{s}). The write pointer
advances one slot per sample and wraps at the end, overwriting the oldest value in
constant time. After $H=128$ new samples arrive (the freshest \SI{2}{s}, teal), the
firmware reads the full ring as the current window and emits a decision; the older
half (navy) is reused as the overlap of the next window.}
\label{fig:ring}
\end{figure}

\subsection{The cost of windowing: dilution at the edges}

Windowing is not free of blind spots, and honesty about them is part of the design.
Because every verdict is a \SI{4}{s} average, a freeze that is \emph{short} relative
to the window, or that is \emph{split} across a window boundary, can be diluted below
the detection threshold and missed.

\begin{itemize}
  \item \textbf{Too short.} The freeze index is $\FI = \text{power}(3\text{--}8)/
        \text{power}(0.5\text{--}3)$ computed over the whole \SI{4}{s} window. If a
        freeze occupies only, say, \SI{1}{s} of that window, the trembling energy is
        averaged together with \SI{3}{s} of calmer signal, pulling the ratio down ---
        possibly below the deployed trip line of $1.815$.
  \item \textbf{Split across a boundary.} A freeze straddling the seam between two
        windows contributes part of its energy to each. Neither window contains the
        whole episode, so each sees a watered-down version, and both can fall short.
\end{itemize}

\begin{figure}[H]
\centering
\resizebox{0.95\textwidth}{!}{%
\begin{tikzpicture}[font=\small]
  \def\sx{0.95}
  % timeline
  \draw[->,thick,ink] (0,0) -- ({13.2*\sx},0) node[right]{\,$t$\,(s)};
  \foreach \t in {0,2,4,6,8,10,12}{
    \draw[ink] ({\t*\sx},0.05) -- ({\t*\sx},-0.05);
    \node[below,ink,font=\footnotesize] at ({\t*\sx},-0.08){\t};
  }
  % --- case A: well-captured freeze (top) ---
  \node[ink,font=\footnotesize,anchor=west] at ({0.1*\sx},3.55)
    {\textbf{Captured:} freeze fills a window $\Rightarrow$ FI high};
  \filldraw[fill=palefreeze,draw=freeze,line width=0.6pt]
    ({2*\sx},2.9) rectangle ({6*\sx},3.35);
  \node[freeze,font=\footnotesize] at ({4*\sx},3.125) {freeze ($\sim$\SI{4}{s})};
  \filldraw[fill=paleteal,draw=teal,line width=0.6pt,rounded corners=1pt]
    ({2*\sx},2.45) rectangle ({6*\sx},2.78);
  \node[teal,font=\scriptsize] at ({4*\sx},2.615) {window aligns};
  \node[grn,font=\footnotesize] at ({8.4*\sx},2.78) {$\FI \uparrow$ detected};
  % --- case B: too-short freeze (middle) ---
  \node[ink,font=\footnotesize,anchor=west] at ({0.1*\sx},1.85)
    {\textbf{Diluted (too short):} \SI{1}{s} freeze in a \SI{4}{s} window};
  \filldraw[fill=palefreeze,draw=freeze,line width=0.6pt]
    ({4.5*\sx},1.2) rectangle ({5.5*\sx},1.62);
  \node[freeze,font=\scriptsize] at ({5*\sx},0.98) {\SI{1}{s}};
  \filldraw[fill=mist,draw=mute,line width=0.6pt,rounded corners=1pt]
    ({3*\sx},0.62) rectangle ({7*\sx},0.93);
  \node[mute,font=\scriptsize] at ({4*\sx},0.775) {3 s calm averaged in};
  \node[freeze,font=\footnotesize] at ({9.0*\sx},0.93) {$\FI$ pulled below trip};
  % --- case C: split across boundary (bottom) ---
  \node[ink,font=\footnotesize,anchor=west] at ({0.1*\sx},-0.6)
    {\textbf{Diluted (split):} freeze straddles two windows};
  \filldraw[fill=palefreeze,draw=freeze,line width=0.6pt]
    ({4*\sx},-1.25) rectangle ({8*\sx},-0.83);
  \node[freeze,font=\scriptsize] at ({6*\sx},-1.04) {freeze across the seam};
  \filldraw[fill=paleteal,draw=teal,line width=0.6pt,rounded corners=1pt]
    ({2*\sx},-1.88) rectangle ({6*\sx},-1.57);
  \node[teal,font=\scriptsize] at ({3.4*\sx},-1.725) {window $k$};
  \filldraw[fill=palenavy,draw=navy,line width=0.6pt,rounded corners=1pt]
    ({6*\sx},-2.30) rectangle ({10*\sx},-1.99);
  \node[navy,font=\scriptsize] at ({8.6*\sx},-2.145) {window $k{+}1$};
  \draw[freeze,dashed,line width=0.7pt] ({6*\sx},-0.78) -- ({6*\sx},-2.34);
  \node[freeze,font=\scriptsize,fill=white,inner sep=1pt] at ({6*\sx},-2.55) {seam};
\end{tikzpicture}}
\caption{When windowing dilutes a freeze. \emph{Top:} a freeze that fills a window
drives the freeze index well above the trip line and is caught. \emph{Middle:} a
sub-window (\SI{1}{s}) freeze is averaged together with \SI{3}{s} of calm signal,
pulling the ratio down. \emph{Bottom:} a freeze split across the seam between two
windows donates half its energy to each, so neither window sees the full episode.}
\label{fig:dilution}
\end{figure}

The \SI{2}{s} hop is the first line of defence here: a freeze missed by one window's
alignment is examined again, \SI{2}{s} later, by the next window --- whose seam falls
in a different place. Overlap thus gives every episode multiple looks. The second
line of defence is downstream: the deployed trip line is biased low ($1.815$ rather
than the Youden-optimal $2.10$), precisely so that a partially-diluted freeze still
clears the bar. The clinical logic is asymmetric --- a missed freeze can mean a fall,
whereas a false alarm merely lights an LED.

\begin{takeaway}
\SI{4}{s} windows give the spectral resolution to separate the \SIrange{3}{8}{Hz}
freeze band; a \SI{2}{s} half-window hop gives a fresh verdict twice a second and a
second look at any episode a single window might dilute. Long enough to see the
freeze, fast enough to report it.
\end{takeaway}
